\textbf{结论名称：}向量加法的平行四边形法则与三角形法则

\textbf{适用条件：}三角形法则对任意两个向量均适用；平行四边形法则要求两个向量不共线，且需将向量平移至同一起点。

\textbf{变量说明：}设向量 $\vec{a} = \overrightarrow{OA}$，$\vec{b} = \overrightarrow{OB}$。三角形法则中将 $\vec{b}$ 平移为 $\overrightarrow{AB}$；平行四边形法则中以 $OA$、$OB$ 为邻边作平行四边形 $OACB$。

\textbf{核心公式：}
\[
\boxed{\overrightarrow{OA} + \overrightarrow{AB} = \overrightarrow{OB}}\quad\text{(三角形法则)}\qquad
\boxed{\overrightarrow{OA} + \overrightarrow{OB} = \overrightarrow{OC}}\ \text{(平行四边形法则，其中$OACB$为平行四边形)}
\]

\textbf{使用场景：}几何向量加法、力的合成、位移合成等矢量合成问题，是向量坐标运算的几何基础。

\textbf{简单例子：}设 $\vec{a}$ 表示向东位移 $3$ m，$\vec{b}$ 表示向北位移 $4$ m。按三角形法则，从 $O$ 向东走 $3$ m 至 $A$，再从 $A$ 向北走 $4$ m 至 $B$，合位移 $\overrightarrow{OB} = \vec{a} + \vec{b}$，长度为 $5$ m，方向北偏东 $\arcsin\frac{3}{5}$。按平行四边形法则，以向东、向北位移为邻边作矩形，对角线即合位移，结果相同。